\section{Theory}

\subsection{Modulus}
We say $b$ is invertible modulo 
$N$ if there exists an integer $c$ such that 
$bc = 1 \mod N$. The multiplicative 
inverse of $b$ modulo $N$ is unique and 
usually denoted $b^{-1}$. When $b$ is 
invertible, we define division as 
$[a/b \mod N] = [ab^{-1} \mod N]$. 

With integers $b, N$ where $b \geq 1$ 
and $N > 1$, $b$ is invertible mod $N$ 
iff $\gcd(b, N) = 1$. 

\subsection{Groups}
A group is a set of elements $\mathcal{G}$
with a group operation $\odot$ on two 
elements, which operates on two elements 
and returns another element in the group. 
\begin{equation}
    \odot: \mathcal{G} \times \mathcal{G} \rightarrow \mathcal{G}
\end{equation}
It has the properties of:
\begin{itemize}
    \item Closure:
    \item Identity
    \item Inverse 
    \item Associativity
\end{itemize}
Abelian groups also have the commutative 
property. 

$\mathcal{Z}_N = \left\{0, \dots, N - 1\right\}$
for $N > 1$ is an abelian group with respect 
to addition modulo $N$. 

The order of a group is the number of 
elements in that group: for the above example,
the order is $N$. 

We define $\mathcal{Z}_N^*$ as the set of 
integers mod $N$ that are relatively prime 
to $N$: this is an abelian group under 
multiplication mod $N$. 

Define $\Phi(N) = |\mathcal{Z}_N^*|$ as 
the order of the group $\mathcal{Z}_N^*$
($\Phi(N)$ is Euler's totient formula).
If $N = pq$ with $p$, $q$ distinct primes, 
then $\Phi(N) = (p-1)(q-1)$.

Given a composite integer $N$, the factoring 
problem is to find integers $p, q > 1$ s.t.
$pq=N$. Factoring is a classic hard problem. 
The smallest prime factor is at most $\sqrt{N}$. 
Consider the following experiment: choose 
two uniform $n$-bit integers $x_1$, $x_2$. 
Set $N = x_1 \times x_2$. Adversary is given 
$N$ and guesses $x_1'$, $x_2'$. Adversary wins 
if $N = x_1' \times x_2'$. This is a 
hard problem if $x_1$, $x_2$ are large primes. 
This is the basis of RSA. 

Let $\mathcal{G}$ be a finite group 
of order $m$. Consider the set $<g> = \{g^0, g^1, \dots\}$.
We saw that $g^m = 1$. Ergo $<g>$ repeats 
after $i$ terms $<g> = \{g^0, \dots, g^{i - 1}\}$. 
If there is an element $g \in \mathcal{G}$ with 
$i = m$, then $<g> = \mathcal{G}$. 
We call $\mathcal{G}$ a cyclic group and $g$ a 
generator of $\mathcal{G}$. This means 
that $\forall h \in \mathcal{G} \exists x \in \{0, \dots, m - 1\}$
s.t. $h = g^x$. For a cyclic group of 
order $m$ with generator $g$ we have 
$\mathcal{G} = <g> = \{g^0, g^1, \dots, g^{m-1}\}$.
We call $x$ the discrete logarithm of $h$ 
w.r.t. $g$ and write $x = \log_g h$. 
For certain groups it's hard to compute 
the discrete logarithm. This is the basis 
of Diffie-Hellman problems. 